% ============================================================================
% Chapter 7: Ancillary Sections
% ============================================================================

\appendix

\tersoo{IMPORTANT NOTE: THE FIRST ANCILLIARY SECTION SHOULD GO OVER ALL RELEVANT POKEMON/SPEEDRUNNING TERMS RELEVANT TO THIS PAPER FOR THE UNINITIATED}

\section{Environment and System Details}\label{sec:appendix_environment}

This appendix provides the technical details underlying the speedrunning benchmark described in Chapter~\ref{sec:benchmark}. We describe the emulator and infrastructure stack, the state extraction pipeline and the role of porymap decompilation data. We also provide a detailed walkthrough of both the tutorial and long-horizon evaluation settings for those less familiar with the storyline of \pokemon{} Emerald.

\subsection{Emulator and Infrastructure}\label{sec:appendix_infrastructure}

The benchmark runs Pok\'{e}mon Emerald on an mGBA emulator~\citep{huderlem2025porymap} operating server-side within a multiprocess architecture. A FastAPI application hosts the primary game server, exposing endpoints for agent interaction through the Model Context Protocol (MCP). The agent communicates with the server over HTTP, issuing tool calls (e.g., \texttt{get\_game\_state}, \texttt{press\_buttons}, \texttt{navigate\_to}) and receiving structured responses containing the current observation.

The emulator runs at a configurable frame rate (default: 80 FPS for gameplay, reduced during heavy computation). Visual frames are captured directly from the mGBA video buffer as $240 \times 160$ pixel images---the native Game Boy Advance resolution. A lightweight frame server streams these frames to a browser-based visualization interface via WebSocket, enabling the real-time evaluation display described in Section~\ref{sec:our_perspective}.

The infrastructure maintains a comprehensive cumulative metric log that records, for every agent step: wall-clock timestamps, milestone completion events, the action taken (button presses and tool calls), token counts for both input and output, API cost, and the full agent reasoning trace (chain-of-thought output). This log serves multiple purposes: it provides the raw data for collection of a wide berth of quantitative metrics beyond those defined in Section~\ref{sec:quantitative_metrics}, enables post-hoc qualitative analysis of agent behavior, and supports automated analysis pipelines where LLMs process trajectory logs to identify failure and success modes within experiments.

The server architecture is designed for extensibility in harness design. New MCP tools can be registered without modifying the core game loop. This is especially important in the interest of extendability, as one can define their own custom-harness by developing new MCP tools and modifying the set of tools the agent has access to. The observation pipeline is modular: additional state channels (e.g., audio, screen regions) can be added by extending the memory reader. The entire system is open-source, enabling participants and researchers to reproduce experiments, develop custom harnesses, and contribute new evaluation scenarios.

\subsection{State Extraction and Porymap Supplementation}\label{sec:appendix_state_extraction}

The structured text state that agents receive (Section~\ref{sec:agent_perspective}) is assembled from two complementary data sources: real-time memory reading and pre-parsed map data from the \texttt{pokeemerald} decompilation project.

\paragraph{Memory Reading.} The memory reader directly accesses GBA RAM addresses to extract game state without relying on visual parsing. Key data includes: the player's position coordinates (from the player object at \texttt{0x03005D8C}), the current map bank and number (resolved to a human-readable location name via an enumeration table), the player's Pok\'{e}mon party (6 slots of 100-byte structures starting at \texttt{0x020244EC}, each containing species, level, HP, status conditions, moves, and PP), the bag inventory (organized into items, key items, Pok\'{e} Balls, TMs, and berries), the game state register (distinguishing overworld, dialog, battle, and menu states), badge flags, and access to a number of variables like the amount of money the player has or the number of pokemon recorded in their Pok\'{e}dex. Memory reading is most important when reasoning about what game state the agent should have access to or for defining milestone completion flags from ground-truth.

\paragraph{Porymap Supplementation.} Raw memory reading provides tile-level map data (metatile IDs, collision flags, elevation), but this information is difficult for an LLM to interpret without semantic context. We supplement memory-read map data with pre-parsed layouts from the \texttt{pokeemerald} decompilation project, processed through Porymap~\citep{huderlem2025porymap}. This supplementation provides simple access to several critical enrichments: NPC positions and types (trainers, shopkeepers, quest-givers), object event locations (items, hidden items, interactive objects), warp coordinates with their destination map and coordinates (enabling the agent to know where doors, stairs, and cave entrances lead), and metatile behavior annotations that classify tiles by function (walkable, surfable, climbable, etc.).

The impact on agent performance is direct: without porymap data, agents must infer spatial structure purely from visual observation and trial-and-error navigation. With porymap data, the agent receives an ASCII map grid with semantically meaningful symbols (\texttt{.}~for walkable, \texttt{\#}~for walls, \texttt{D}~for doors, \texttt{N}~for NPCs, \texttt{I}~for the player) and a list of known portal coordinates with destination names. Importantly, by providing a global-view of the current map location, beyond what the agent can see in the visual frame, the inclusion of Porymap data represents a fundamental aspect of our Environment Design which keeps the process of navigation challenging, but doable. This data is also a prerequisite for the pathfinding tool we implement that serves as one of the most valuable and frequently called MCP tools in our expert harness \tersoo{NOTE: we should be consistent in when we use H-auto, H-expert, H-min vs 'expert harness' 'minimal harness etc.}.

\subsection{Walkthrough: Tutorial (Littleroot to Roxanne)}\label{sec:appendix_tutorial_walkthrough}

The tutorial evaluation spans 15 milestones \tersoo{NOTE: this list doesnt track the exact 15 milestones perfectly, so i removed the enumeration. Please reference the figure with the tutorial milestones}, each testing a distinct combination of agent competencies. We describe the progression and the capabilities each segment exercises, referencing the open Pok\'{e}mon Emerald walkthrough~\citep{bulbapedia2025walkthrough} as a canonical game guide.

\textbf{Littleroot Town} (Start). The agent begins in a moving truck outside of their house in Littleroot Town. The specific location of the house is dependent on their selected gender, though our benchmark defaults to the male character, "Brendan" for all evaluations. The player must then walk into their bedroom, set their clock, before navigating downstairs, talking to their mother, exiting the house, and walk to the neighboring house.

\textbf{Route 101 / Rescue Prof.\ Birch.} The agent encounters Professor Birch being chased by a wild Pok\'{e}mon and must select a starter Pok\'{e}mon (Treecko, Torchic, or Mudkip) from his bag. The choice of starter strongly correlated with their completion time of the first gym as Torchic is week to the rock typing of the first gym, while Treecko and Mudkip learn supereffective moves against rock types.

\textbf{Oldale Town.} A brief traversal milestone through the first town, north of Littleroot town with a Pok\'{e}mon Center and shop.

\textbf{Route 103 / Rival Battle.} The agent must navigate north of Oldale town, through grass that may trigger encounters with wild pokemon and battle the rival trainer. This represents the first scripted trainer battle, testing basic battle strategy against an opponent who always has a starter with a more favourable type combination against the player.

\textbf{Return to Birch's Lab.} The agent must backtrack to Littleroot Town from Oldale town, through Route 101, and into Professor Birch's Lab in Littleroot town to receive the Pok\'{e}dex and Pok\'{e} Balls. 

\textbf{Route 102.} The route immediately west of Oldale town chalk full of wild encounters and several mandatory trainer battles. Tests managing wild encounters while progressing.

\textbf{Petalburg City / Dad's Gym.} The agent meets the player's father (Norman, the gym leader) but cannot challenge him yet. They must navigate back to Route 102, where they must help Wally catch a pokemon. This tests dialog dialogue handling and comprehension.

\textbf{Route 104 (South).} Continuation westward of Petalburg city with optional trainers, items, and wild encounters.

\textbf{Petalburg Woods.} A maze-like forest area with several trainers, an abundance of wild grass, and a mandatory encounter against a Team Aqua Grunt. This is the first ``dungeon'' environment testing navigation in complex terrain with a required scripted battle at a distance that makes backtracking to heal the party expensive.

\textbf{Route 104 (North).} Emergence from the woods and passage through the northern section with wild grass for optional wild encounters and a number of optional and mandatory trainer battles.

\textbf{Rustboro City.} Arrival at the destination city from Route 104 (North), harboring the Rustboro Gym

\textbf{Rustboro Gym / Gym Trainers.} The gym contains trainers that must be defeated before challenging the leader. This tests sequential battle management, puzzle navigation through a simple maze, and resource conservation.

\textbf{Defeating Roxanne.} The gym leader battle against Roxanne and her Rock-type team consisting of several Geodudes and a powerful level 15 Nosepass. This split tests strategic battle planning, especially against a more sophisticated adversary with the ability to heal their \pokemon{}. Treecko and Mudkip starters have a distinct type advantages, while Torchic requires more careful play either through the acquisition of pokemon with a more favourable type combination or through evolution of Torchic to Combusken, to learn fighting-type moves which are super-effective against Rock-types. Once the player defeats Roxanne.
\end{enumerate}

\subsection{Walkthrough: Long Horizon (Through Gym 3)}\label{sec:appendix_longhorz_walkthrough}

After the tutorial, progression becomes substantially more complex. We outline the major segments through Gym~3, emphasizing the structural features that make this setting challenging for autonomous agents.

\paragraph{Rustboro City $\to$ Rusturf Tunnel.} After defeating Roxanne, the agent must pursue a Team Aqua Grunt who stole the Devon Goods from the researcher earlier in Petalburg Woods. This requires navigating east to Route~116 where more wild encounters and trainers await and into Rusturf Tunnel, where the grunt is confronted and the goods recovered. This segment formally introduces the side-quest dependency pattern: certain segments of the main gym quest cannot proceed until a corresponding side-quest narrative thread is resolved.

\paragraph{Devon Corporation.} The agent returns the Devon Goods to the Devon Corporation president in Rustboro City, receiving a letter to deliver to Steven Stone in Dewford Town and a package for Captain Stern in Slateport City. This fetch quest creates an additional branching dependency: the agent now has two objectives in different locations (Dewford Gym completion vs. Granite Cave navigation), and the order of completion affects routing efficiency.

\paragraph{Mr.\ Briney's Cottage $\to$ Dewford Town.} The agent must travel south to Route~104 and locate Mr.\ Briney's cottage, where the old sailor agrees to ferry the player to Dewford Town. This imposes a non-trivial navigation requirement as the cottage is in a previously visited area, requiring the player to recognize the need for backtracking.

\paragraph{Gym 2: Brawly (Dewford Town).} Dewford's gym features a dark room that lights up incrementally as the player defeats trainers a spatial puzzle requiring the agent to navigate partially blind terrain. Brawly's Fighting-type team tests battle strategy with different type matchups than Roxanne, and upon succesful completion, the player receives the knuckle badge. Note that the Knuckle Badge and the Dynamo badge can be completed in either order.

\paragraph{Granite Cave / Steven Delivery.} The agent must explore Granite Cave to find Steven Stone and deliver the letter. This multi-floor cave requires the HM move Flash (optional but significantly helpful in navigation), introducing the HM dependency system where progress requires teaching specific moves to Pok\'{e}mon. It is important to note that this delivery gates the players ability to access Slateport City via Mr. Briney's ferry, a prerequisite for navigating to Mauville City to challenge for the third gym. 

\paragraph{Sailing to Slateport City.} Mr.\ Briney ferries the agent to Slateport City, where Captain Stern must be found at the Oceanic Museum. This segment involves a second Team Aqua encounter, further advancing the interleaved narrative arc.

\paragraph{Route 110 $\to$ Mauville City.} The northward route to Mauville City features the Cycling Road (requiring a Bicycle, obtained in Mauville, though not strictly necessary for split completion) and a second rival battle at the route's midpoint. The route is moderately long with multiple trainers.

\paragraph{Gym 3: Wattson (Mauville City).} Wattson's gym features an electric fence puzzle where the agent must navigate a grid of switches to open the path to the gym leader. This puzzle is non-trivial and requires careful navigation to certain switches and around barrier positions that change with the activation of switches. His Electric-type team (Voltorb, Electrike, Magneton, Manectric) poses a significant challenge, especially for agents whose teams lack Ground-type coverage. This gym represents a meaningful difficulty spike and serves as the endpoint of the long-horizon evaluation.

\vspace{1em}
The progression from Gym~1 through Gym~3 illustrates how the game's complexity compounds: each segment introduces new mechanics (HMs, ferry travel, dark rooms, electric puzzles), new narrative dependencies, new resource management pressures, and a set of optional story elements that the agent must often reason about skipping in service of an efficient completion time. An agent that reaches Gym~3 has demonstrated not only the basic competencies tested in the tutorial but also the ability to maintain coherent long-term plans, manage an evolving team, and navigate a web of interconnected objectives across many hours of autonomous play. We find that success in this split is highly predictive of continued progress through even longer stretches of the game that may take multiple days of contiguous execution.

\section{Competition Participant Methodologies}\label{sec:appendix_competitor_methods}

This section presents methodology summaries contributed by top-performing teams in the NeurIPS 2025 \pokeagent{} Challenge Speedrunning Track ~\citep{karten2025pokeagent}.

\subsection{Heatz (Speedrunning Track Champion)}\label{sec:heatz}

\textbf{Author:} Junik Bae

\begin{figure}[h]
  \centering
  \includegraphics[width=\textwidth]{appendix_competitors/heatz_figure.pdf}
  \caption{\textbf{Overview of Scripted Policy Distillation (SPD).} The approach consists of three stages: (1) Subgoal Generation where an LLM decomposes the task into sequential subgoals with executable success conditions; (2) Scripted Policy Generation where the LLM generates policies that can invoke VLM tools and use self-directed logging for debugging; (3) Script-Guided RL where scripted policies are distilled into neural networks via supervised learning followed by RL with expert action guidance.}
  \label{fig:heatz_method}
\end{figure}

Reinforcement learning offers a powerful framework for training autonomous agents, yet applying RL directly to long-horizon tasks with sparse rewards remains fundamentally challenging. To address this, Heatz proposed leveraging LLM-generated scripted policies as priors for RL exploration (Figure~\ref{fig:heatz_method}). Although these scripted policies are not always optimal, they facilitate RL training by initializing exploration from a distribution that already reaches the goal, rather than from scratch. The framework, \textbf{Scripted Policy Distillation (SPD)}, consists of three stages: (1) subgoal generation, (2) scripted policy generation, and (3) script-guided RL.

\paragraph{Subgoal generation.} Given a long-horizon task specification, an LLM decomposes the task into a sequence of subgoals, each paired with an executable success-condition function \texttt{success\_cond(state)} that determines subgoal completion.

\paragraph{Scripted policy generation.} For each subgoal, the LLM generates a scripted policy that maps states to actions. The scripted policy interacts with the environment until \texttt{success\_cond(state)} returns \texttt{True} or a timeout occurs. Upon failure, the LLM analyzes execution traces and revises either the policy code or the subgoal specification. Policies can optionally query a vision-language model (Qwen3-VL-8B) for visual cues, and self-directed logging provides concise summaries of agent interactions rather than full execution trajectories.

\paragraph{Script-guided RL.} Once all scripted policies reliably achieve their subgoals, they are distilled into neural network policies via imitation learning on expert trajectories, followed by RL with expert action guidance. A DQN agent is trained with expert trajectories seeded into the replay buffer and expert actions executed with annealed probability during rollouts.

SPD achieved a 40:13 completion time. The resulting policy exhibited emergent behaviors: BFS-based local pathfinding during scripted policy generation, and after distillation, shorter routes, trainer battle skipping, and faster menu interactions discovered through RL optimization.

\subsection{Hamburg Pok\'{e}Runners (Second Place)}\label{sec:hamburg}

\textbf{Team:} Benedikt Schink, Arian Urdu, Matin Urdu

\begin{figure}[h]
    \centering
    \includegraphics[width=0.8\textwidth]{appendix_competitors/hamburg_architecture.pdf}
    \caption{Hamburg Pok\'{e}Runners' recurrent PPO architecture, showing the CNN encoder for visual input, MLP for positional and milestone encoding, and LSTM for temporal reasoning.}
    \label{fig:hamburg_architecture}
\end{figure}

Hamburg Pok\'{e}Runners achieved second place (first badge: 01:14:43) using a reinforcement learning approach based on recurrent proximal policy optimization (PPO). The architecture consists of a CNN encoder receiving downsampled game frames (grayscale, $128 \times 128$ pixels), sinusoidal positional encodings for in-game coordinates, one-hot location encoding, and a binary milestone vector $m \in \{0, 1\}^{38}$ encoding sub-milestones and visited locations. An LSTM in the recurrent state provides temporal reasoning (Figure~\ref{fig:hamburg_architecture}).

The milestone vector played a crucial role in stability: it served both as a memory component (tracking achievements) and as goal conditioning (tracking the current objective), making the model more robust to the inherent randomness of Pok\'{e}mon Emerald and non-deterministic input/output latency. The team observed simple error-correcting behavior where the agent would overshoot or undershoot targets and backtrack to correct its trajectory.

The reward function includes only positive rewards: competition sub-milestones, locations visited, coordinate exploration, badges, money earned, HP healed, Pok\'{e}mon leveled, and damaging attacks used. All negative rewards (e.g., step penalties) caused model degradation as the agent reverted to safe, non-progressive behavior. Due to limited computational resources, multiple smaller models were trained and stitched together during inference.

\subsection{Anthonys (Third Place)}\label{sec:anthonys}

\textbf{Author:} Anthony Sistilli

Anthonys' approach centered on decomposing game progression into discrete phases with specialized navigation strategies for each context. The core innovation was combining deterministic pathfinding algorithms with context-aware prompt engineering.

\textbf{Phase-Based State Management}: The challenge was structured into seven distinct phases corresponding to major game milestones. Each phase maintained its own prompt template with conditional instructions that dynamically adapted based on completed objectives and current location.

\textbf{A* Pathfinding with Directional Priorities}: Rather than relying solely on the LLM for spatial reasoning, A* pathfinding was implemented on the game map grid to compute optimal movement sequences. The system extracted walkable tiles from game memory, constructed a navigable grid, and generated action sequences to reach target coordinates.

\textbf{Battle State Management}: Battle sequences required special handling due to the game's menu system complexity. An input-clearing mechanism sent predetermined button sequences to ensure the agent escaped bad states.

\textbf{Adaptive Recovery and Stuck Detection}: The system tracked position history to detect when the agent remained stationary for multiple turns, indicating an obstacle or navigation failure.

\subsection{Deepest (Judge's Choice: Most Efficient)}\label{sec:deepest}

\textbf{Team:} Seonghun Hong, Hyunyoung Jeong, Hyeokgi Kim, Jaeyoon Jung \\
\textbf{Affiliation:} Seoul National University, Soongsil University, MAUM.AI

\begin{figure}[h]
    \centering
    \includegraphics[width=0.92\textwidth]{appendix_competitors/deepest_architecture.png}
    \caption{Deepest agent architecture. The agent receives partial observations and milestone-specific guidance from the Guidebook. Gemini 2.5 Flash reasons and outputs actions, optionally invoking tools (Pathfinding, Memory, Thinking). Actions execute in the emulator, producing the next observation.}
    \label{fig:deepest_architecture}
\end{figure}

Deepest proposed a training-free, fully autonomous agent architecture that operates under human-like perceptual constraints using Gemini 2.5 Flash. The approach achieved 5th place by time (02:04:29) but received the Judge's Award for most sample-efficient completion, completing in just 649 actions---the fewest of any team (Figure~\ref{fig:deepest_architecture}).

A core design principle is the absence of privileged game state: the agent receives no ground-truth map information, and all navigation decisions are computed exclusively from partially observed tiles explored during gameplay. To support high-level strategy, a \textbf{Guidebook} system provides milestone-specific knowledge (e.g., selecting Mudkip for type advantage). \textbf{Image goals}---reference images of destination locations---enable the agent to visually recognize objectives.

Three auxiliary tools augment the agent's capabilities:
\begin{itemize}[nosep]
    \item \textbf{Pathfinding Tool}: coordinate-based navigation for known destinations and directional exploration for unknown ones, using A* on the observed map grid with game-specific constraints.
    \item \textbf{Memory Tool}: explicit persistence of critical state information with a sliding window policy retaining memory over the most recent 10 conversation turns.
    \item \textbf{Thinking Tool}: adaptive computation toggling the model's reasoning budget---2048 tokens for complex decisions, 512 for routine navigation.
\end{itemize}

\section{Expert Harness Tool Reference}\label{sec:appendix_tools}

This section provides a complete reference for the MCP tools and subagent tool access configurations in $\mathcal{H}_{\text{expert}}$ as described in Section~\ref{sec:system_design}.

\subsection{Tool Access by Agent Role}\label{sec:appendix_tool_access}

Table~\ref{tab:tool_access} details which tools are available to each agent role within the expert harness. The orchestrator has the broadest access; subagents receive restricted tool sets to enforce separation of concerns. Tools marked with $\dagger$ are available only in $\mathcal{H}_{\text{expert}}$ (stripped in $\mathcal{H}_{\min}$ and $\mathcal{H}_{\text{auto}}$). Custom subagents receive a configurable tool set specified at creation time, subject to the constraint that \texttt{execute\_custom\_subagent} is always forbidden (no recursive subagent spawning).

\begin{table}[h]
\centering
\caption{Tool access matrix by agent role. \cmark{} = available, \xmark{} = unavailable, $\dagger$ = $\mathcal{H}_{\text{expert}}$ only.}
\label{tab:tool_access}
\small
\begin{tabular}{l c c c c}
\toprule
\textbf{Tool} & \textbf{Orchestrator} & \textbf{Battler} & \textbf{Planner} & \textbf{1-Step} \\
\midrule
\multicolumn{5}{l}{\textit{Game Control}} \\
\texttt{press\_buttons}           & \cmark & \cmark & \xmark & \xmark \\
\texttt{complete\_direct\_objective} & \cmark & \xmark & \xmark & \xmark \\
\texttt{navigate\_to}$^\dagger$   & \cmark & \xmark & \xmark & \xmark \\
\midrule
\multicolumn{5}{l}{\textit{Memory}} \\
\texttt{process\_memory}          & \cmark & \cmark & \cmark & \xmark \\
\texttt{add\_memory}              & \xmark & \cmark & \cmark & \xmark \\
\texttt{search\_memory}           & \xmark & \cmark & \cmark & \xmark \\
\texttt{get\_memory\_summary}     & \xmark & \cmark & \cmark & \xmark \\
\midrule
\multicolumn{5}{l}{\textit{Skills \& Code Execution}} \\
\texttt{process\_skill}           & \cmark & \cmark & \cmark & \xmark \\
\texttt{run\_skill}               & \cmark & \xmark & \xmark & \xmark \\
\texttt{run\_code}                & \cmark & \xmark & \xmark & \xmark \\
\midrule
\multicolumn{5}{l}{\textit{Research \& Planning}} \\
\texttt{get\_progress\_summary}$^\dagger$ & \cmark & \cmark & \cmark & \xmark \\
\texttt{get\_walkthrough}$^\dagger$ & \cmark & \xmark & \cmark & \xmark \\
\texttt{lookup\_pokemon\_info}$^\dagger$ & \cmark & \cmark & \cmark & \xmark \\
\texttt{replan\_objectives}       & \xmark & \xmark & \cmark & \xmark \\
\midrule
\multicolumn{5}{l}{\textit{Subagent Management}} \\
\texttt{process\_subagent}        & \cmark & \xmark & \xmark & \xmark \\
\texttt{execute\_custom\_subagent} & \cmark & \xmark & \xmark & \xmark \\
\bottomrule
\end{tabular}
\vspace{0.5em}

\footnotesize{1-step subagents (\texttt{reflect}, \texttt{verify}, \texttt{summarize}, \texttt{gym\_puzzle}) receive no tool access; they perform a single VLM pass over a trajectory window segment. The planner additionally has access to the four built-in 1-step subagents and \texttt{replan\_objectives} for objective CRUD operations.}
\end{table}

\subsection{MCP Tool Definitions}\label{sec:appendix_tool_defs}

Table~\ref{tab:tool_defs} provides brief descriptions of all MCP tools exposed by the game server, plus locally defined tools in the agent process.

\begin{table}[h]
\centering
\caption{Complete MCP tool reference for the \pokeagent{} harness.}
\label{tab:tool_defs}
\small
\begin{tabular}{l p{9.5cm}}
\toprule
\textbf{Tool} & \textbf{Description} \\
\midrule
\multicolumn{2}{l}{\textit{Game Control}} \\
\texttt{get\_game\_state}        & Returns current observation: screenshot, ASCII map, party, bag, location, coordinates, and game context (overworld/battle/dialog/menu). Called automatically each step. \\
\texttt{press\_buttons}          & Executes a sequence of GBA button presses (\texttt{A}, \texttt{B}, directional, \texttt{START}, \texttt{SELECT}, \texttt{L}, \texttt{R}, \texttt{WAIT}). \\
\texttt{complete\_direct\_objective} & Marks the current direct objective as completed and advances to the next in the specified category (story, battling, or dynamics). \\
\texttt{navigate\_to}            & A* pathfinding to target $(x,y)$ coordinates with configurable path variance and NPC collision avoidance. \\
\midrule
\multicolumn{2}{l}{\textit{Persistent Memory (CRUD)}} \\
\texttt{process\_memory}         & Unified CRUD for hierarchical long-term memory entries (read, add, update, delete). \\
\texttt{add\_memory}             & Add a single entry to the long-term memory store (legacy interface). \\
\texttt{search\_memory}          & Keyword search over long-term memory entries. \\
\texttt{get\_memory\_summary}    & Retrieve a summary of the long-term memory store. \\
\texttt{get\_memory\_overview}   & Compact tree view of all memory entry IDs and titles. \\
\midrule
\multicolumn{2}{l}{\textit{Skill Library}} \\
\texttt{process\_skill}          & Unified CRUD for skill library entries (behavioral or executable). \\
\texttt{get\_skill\_overview}    & Compact tree view of all skill IDs and names. \\
\texttt{run\_skill}              & Execute a skill's Python code in the sandbox with game tool access. \\
\texttt{run\_code}               & Execute arbitrary Python in the sandbox (read-only game tools) for prototyping. \\
\midrule
\multicolumn{2}{l}{\textit{Subagent Registry}} \\
\texttt{process\_subagent}       & Unified CRUD for the subagent registry (create, inspect, update, delete entries). \\
\texttt{get\_subagent\_overview} & Compact tree view of all subagent IDs and names. \\
\midrule
\multicolumn{2}{l}{\textit{Research \& Progress}} \\
\texttt{get\_progress\_summary}  & Milestones, location, objective status, completed objectives, and memory overview. \\
\texttt{get\_walkthrough}        & Retrieve a section of the Pok\'{e}mon Emerald walkthrough (Parts 1--21). \\
\texttt{lookup\_pokemon\_info}   & Query Bulbapedia for Pok\'{e}mon stats, moves, evolution, and location data. \\
\texttt{list\_wiki\_sources}     & List available wiki data sources for \texttt{lookup\_pokemon\_info}. \\
\midrule
\multicolumn{2}{l}{\textit{Objective Management}} \\
\texttt{get\_full\_objective\_sequence} & Return the complete objective state across all categories. \\
\texttt{replan\_objectives}      & Apply index-based edits (add, complete, reorder, delete) to objectives within a single category. \\
\midrule
\multicolumn{2}{l}{\textit{Spatial Memory (Experimental)}} \\
\texttt{save\_map}               & Save the agent's mental map of a location to persistent storage. \\
\texttt{load\_map}               & Load a previously saved mental map for a location. \\
\bottomrule
\end{tabular}
\end{table}

\section{Example Prompts for Ablation Studies}\label{sec:appendix_prompts}

\textit{[To be written. Representative system prompts, base prompts, and evolution outputs across experimental conditions.]}

\section{Gameplay Demonstrations}\label{sec:appendix_gameplay}

\textit{[To be written. Selected trajectory excerpts illustrating agent behavior across different harness configurations. Annotated reasoning traces.]}

\section{AutoEvolve Store Comparisons}\label{sec:appendix_stores}

\textit{[To be written. Full memory, skill, and subagent store snapshots across AutoEvolve experiments, comparing evolved artifacts across trials and settings.]}

\section{Engineering Standards}\label{sec:appendix_engineering_standards}

Per the ECE Department's Independent Work and Thesis Handbook, this section identifies the engineering standards and broadly accepted industrial conventions used in this thesis.

\textit{[To be completed. The following categories should be addressed with specific standards, tools, and their role in the project:]}

\begin{itemize}
    \item \textbf{Software frameworks and libraries.} \textit{[e.g., Python 3.x, FastAPI, mGBA emulator SDK, Google Generative AI SDK, NumPy. Describe role of each in the system.]}
    \item \textbf{Communication protocols.} \textit{[e.g., HTTP/REST for agent--server communication, WebSocket for frame streaming, Model Context Protocol (MCP) for tool integration. Cite relevant specifications.]}
    \item \textbf{Data formats and serialization.} \textit{[e.g., JSON for structured state representation and tool call/response payloads, PNG for frame capture, BibTeX for bibliography management.]}
    \item \textbf{Evaluation methodology.} \textit{[e.g., wall-clock timing standards, milestone-based progress metrics, token counting conventions, cost estimation methodology.]}
    \item \textbf{Reproducibility and open-source practices.} \textit{[e.g., version control (Git), open-source licensing, public benchmark release, documentation standards.]}
    \item \textbf{Security and sandboxing.} \textit{[e.g., restricted built-in functions in the code execution sandbox, process-local isolation for agent-generated code.]}
    \item \textbf{Game emulation.} \textit{[e.g., mGBA emulator accuracy standards, GBA hardware specifications ($240 \times 160$ resolution, ARM7TDMI processor), ROM integrity verification.]}
\end{itemize}
